\chapter[Justice: Equality or Due Shares?]{Is Justice Giving Everyone the Same, or Giving Each Their Due?}
\section{A Knife and a Cake}
First, pick up an object: an ordinary fruit knife.

Nothing about the knife is special. What matters is what it will do next. It is Grandfather's birthday. An eight-inch cream cake sits on the table, ringed with strawberries of different sizes. There are seven people in the family, and the knife has been handed to you. From this moment it is no longer cutlery but an instrument of power: it determines how sweet everyone's evening will be.

Several problems soon emerge. Your little sister has been staring at the largest strawberry for ten minutes. Your cousin spent an hour massaging Grandfather's back this afternoon and thinks he deserves an extra slice. Your father has high blood sugar, and his doctor has told him not to eat cream. If you give him a slice as large as everyone else's, is that really fair? Then there is you: the person cutting knows best which slice is largest. Can you resist?

A clever popular rule says, "I cut, you choose": the cutter takes the last piece. Introduce that rule and the cutter instantly becomes the world's most precise scale, because any undersized slice will end up on their own plate. Notice how seriously humans have always taken distribution. We have even invented an allocation mechanism that runs automatically without supervision.

But notice the assumption hidden inside it: everyone ought to receive equally large pieces. What if someone wants to distribute them by age, contribution, or eagerness to eat? "I cut, you choose" cannot answer. It executes equal division down to the millimeter without explaining why the division should be equal.

This chapter concerns the place where that knife really cuts: does justice mean everyone receives the same, or each receives their due, the share they ought to have? If the latter, who decides what each person deserves? You encounter this question at dinner and in class meetings. Later, it will return on payslips, hospital registration forms, and court judgments, again and again.

\section{Ten Minutes in a Class Meeting}
Now come with me into a particular scene.

Time: a Friday near the end of May, twenty past four in the afternoon, during a class meeting. Place: Grade Eight, Class Three, at an ordinary secondary school. Ceiling fans creak overhead, unable to shift June's oppressive heat. In the blackboard's upper-right corner is a countdown to the college entrance examination, intended for the senior-year teaching building. This classroom has been borrowed as a reserve examination room, and the desks and chairs have not yet been put back.

The class teacher, Ms. Li, writes on the board: "Ten tablets: how should we distribute them?"

Here is what happened. A local business donated educational tablets to the school. The principal's administrative meeting decided to give each class ten, managed collectively and used primarily for learning. Grade Eight, Class Three has forty-eight students and ten tablets. Ms. Li does not want to decide alone. They will discuss it together this period.

The room falls silent for three seconds, then boils over.

Dachuan raises his hand first. A top mathematics student, he always speaks directly. "I think they should go to the ten students whose marks have improved most. The company donated them for learning, so the best students should use them. That would encourage everyone to study. If everything is divided equally, what's the point of trying?"

Lin Wan sits by the window in the third row. Her voice is quiet, but many hear her. "My family... doesn't have a computer. I watch online lessons on a phone. The screen's so small it hurts my eyes." She does not say, "Give one to me." She sits down, but everyone understands the unspoken ending.

Xiaoman jumps up impatiently. "But you can't distribute them by marks! Good students already have better conditions. Extra classes cost tens of thousands of yuan a year. I say draw lots! No favorites, no complaints. Leaving it to chance is fairest."

Near the back, Chen Guo gives a cold laugh. He belongs to the school's athletics team and trains until half past six after school. "Whatever you do, I won't get one. I can't beat Dachuan's marks, and I can't bring myself to compete over hardship. I have one question: why should whoever talks fastest get to decide?"

Within ten minutes, four proposals appear on the board: reward marks, prioritize need, draw lots for everyone, or take turns. Ms. Li crosses none out. She simply adds a smaller question: "Any others?" The argument grows louder. Someone at the back starts drumming a pencil case against a desk.

Remember this classroom. The rest of the chapter answers the questions that erupted in this lesson.

\section{When Sameness and Desert Clash}
Before judging who is right, translate the argument.

On the surface, four people are arguing over ten tablets. Beneath that, each is speaking a completely different word. We usually throw these words into one large basket called fairness. Once there is something to distribute, they run in separate directions.

\textbf{Equality}: everyone receives the same share. Xiaoman's lottery is a variant: if the shares cannot be equal, make everyone's chance of receiving one equal. Cutting a cake into eight equal slices follows this logic.

\textbf{Need}: those lacking most come first. This is Lin Wan's reason. Emergency departments triage by medical urgency rather than registration order. If someone with a nosebleed arrives alongside someone having a heart attack, nobody thinks first come, first served makes sense here.

\textbf{Contribution}: those who put in more or perform better receive more. This is Dachuan's reason, the logic of piece rates, pay according to work, and scholarships.

\textbf{Rights}: some things never enter the distribution debate because everyone is entitled to them from birth, such as education or freedom from arbitrary beating and verbal abuse. Chen Guo's question about why fast talkers should decide touches this boundary. He suspects that the discussion's procedure itself violates something everyone should have.

Then there is \textbf{fairness}, the most troublesome word of all. Often it is not a sixth standard but an overall judgment about whether the preceding standards should give way in a particular situation. When people ask whether something is fair, they are usually arguing about which standard applies here.

Now the real question can reach the table. Actually, there are two, and they press against each other.

First: does justice mean treating everyone alike, or giving each their due, treating people differently according to their circumstances, efforts, and needs?

The second is more uncomfortable. If we choose each person's due, who holds the ruler? With need, who determines who needs most? With contribution, does a student with poor marks who explains problems to classmates every day count as contributing? With equality, is giving your father with high blood sugar the same-sized slice as everyone else a kind of deliberate blindness?

Neither question yields an answer at a glance. Before continuing, take a side in the classroom. If you were Ms. Li, which of the four proposals would you most want to cross out now? Remember your answer and revisit it after reading.

\section{Give Every Side a Full Hearing}
In a noisy classroom argument, it is easy to hear only your own side. Let us do something harder: state each side's reasons as fully as possible, even if you disagree. In debate, this means explaining the other person's case so well that they themselves nod before you rebut it. Defeating a weakened opponent is no real victory.

\textbf{For Dachuan: rewarding excellence is an engine, not favoritism.}

Dachuan's case has three increasingly substantial layers. First, effort needs a response. If studying well and poorly bring the same result, the ten most improved students wasted their late nights. Who will stay up next time? A group that never responds to excellence eventually becomes one in which nobody wants to excel. Second, resources should go where they produce most. A tablet is a learning tool. Giving it to whoever can learn best with it resembles giving seeds to the most skilled farmer. That sounds cold, but it maximizes the class's harvest. Third, rules should be announced in advance and applied consistently. Scholarships work this way; nobody can replace the system with a lottery on payment day. Having heard these three layers seriously, you cannot dismiss them with "You're just selfish."

\textbf{For Lin Wan: with unequal starting points, the same is not the same.}

Lin Wan's central argument is that allocation should attend to real lives, not places on a blackboard. The same tablet becomes a fourth screen in a home with a laptop, iPad, and printer. In her home, it spans the entire distance between being unable and able to see online lessons clearly. Looking only at what each person receives, without asking what it changes, treats people like shelves. Nor did she choose her difficulty: nobody submitted a preference form before being born. Making a child pay for her parents' income is the real unfairness.

\textbf{For Xiaoman: a lottery limits the damage from competing standards.}

Xiaoman's insight is that whenever a standard requires human assessment, the authority to assess is itself power. Who identifies the greatest need? Must family income be displayed? Making Lin Wan stand before the class and describe her family's hardship harms her through the very procedure meant to prioritize need. A lottery abandons assessment and with it the opportunities to exploit assessment power for private advantage. It admits that when competing standards cannot produce a decision, acknowledging "We cannot judge" is more honest than pretending "We can judge accurately."

\textbf{For Chen Guo: procedure is also distribution.}

Chen Guo's voice is easiest to dismiss as grumbling, but his question cuts deepest. Loud voices, early raised hands, and larger circles of friends are already deciding the outcome. Procedure is not a ceremony before distribution; it is part of distribution. What it allocates is the right to speak.

Look elsewhere in the world and this argument turns out to be an old human routine.

On Tanzania's grasslands, Hadza people still live by hunting and gathering. When hunters return with large game, meat is shared throughout the camp. The bow's owner does not keep it all; often even the hunter takes no extra. Anthropologists observing them have noted a strikingly simple reason: today I caught something; tomorrow I may go hungry. Sharing is not saintly morality but a safety net everyone weaves together, because anyone may need it.

In Athens twenty-five hundred years ago, the city-state applied lotteries to politics. Many public offices were filled not by election but by random selection among citizens. Their reason resembled Xiaoman's exactly: elections always favor eloquence, fame, and wealth. Only lots could turn public office from a rich person's privilege into an equal opportunity. Random jury selection in some countries today descends from this old mechanism.

We will read an East Asian version more than two thousand years old in the next section. For now, remember that how to distribute, who distributes, and what happens when distribution goes wrong have never troubled just one place.

\section[Five Questions About Distribution]{Thinking Bridge: Ask Five Questions About Distribution}
Whenever an argument arises about allocation, from classroom cleaning duties to national policy, do not immediately take sides. Bring out a five-question checklist and untangle one knot into five threads. This is the chapter's portable tool.

\textbf{First: what exactly is being distributed?} Cake, tablets, scholarships, opportunities to criticize, rest time, or the right to be heard? Things include more than objects: opportunities, time, and dignity are distributed too. Many arguments arise because both sides think they are dividing the same thing, when one means tablets and the other means dignity.

\textbf{Second: who is distributing it?} Making the cake cutter choose last guards against the fact that allocators have preferences of their own. Asking whether rule-makers are governed by their rules exposes half of unfairness.

\textbf{Third: by which standard?} Equality, need, contribution, rights: each is a different ruler. A mature discussion does not select the only ruler but explains which takes priority here, and why.

\textbf{Fourth: is the procedure just?} Who set the standard? Were affected people present? Are there explanations, public disclosure, and appeals? A good procedure need not produce your preferred outcome. A bad one will eventually produce an outcome nobody can accept.

\textbf{Fifth: what if the allocation is wrong?} How are wrongful judgments reversed? How should damage to another person's belongings be compensated? Should a teacher publicly apologize to a falsely accused student? What happens afterward matters as much as the allocation itself.

Behind these five questions stand three concepts with formal names that you will encounter repeatedly in news and textbooks:

\begin{itemize}
\item \textbf{Distributive justice}: how should goods, opportunities, and burdens be allocated? All four classroom proposals concern it.
\item \textbf{Procedural justice}: is the decision-making process fair? Were rules published beforehand, applied equally, and accompanied by opportunities for those involved to speak? The dispute between fastest registration and drawing lots concerns it.
\item \textbf{Corrective justice}: after harm occurs, how should it be remedied through compensation, apology, exoneration, or punishment? Should Chen Guo pay for someone else's racket he broke? Does anyone address three years of harm to Lin Wan's eyes from a small screen? These questions concern it.
\end{itemize}
Next time a group chat argues about fairness, ask which level is at issue. In eight arguments out of ten, one side discusses distribution while the other discusses procedure, and both think they have won.

\section{Distribution Three Thousand Years Ago}
Now read several short primary sources. The chapter's questions are older than you might imagine, and ancient answers differ from one another. That itself matters.

The first comes from the "Evolution of Rites" chapter of China's Book of Rites, describing an ideal society in words attributed to Confucius:

\begin{quote}
When the Great Way prevailed, the world belonged to all... Thus people cherished not only their own parents and children, but ensured that the old could live out their years, adults could put their abilities to use, and the young could grow; widowers, widows, orphans, elders alone, and those disabled or ill all received support.
\end{quote}
Broadly, when the Great Way operates, the world is shared. People care not only for their own parents and children but ensure that elders end life well, adults have useful work, children grow, and widowers, widows, orphans, solitary elders, and disabled people receive support. Notice the final part: everyone listed is among those in greatest need. This is one of the earliest explicit statements of distribution according to need found anywhere, preceding national welfare systems by more than two thousand years.

The second comes from the Analects, "The Ji Family":

\begin{quote}
I, Qiu, have heard that those who govern states and houses worry not about scarcity but about uneven distribution, not about poverty but about insecurity.
\end{quote}
In broad terms, rulers worry less about insufficient wealth than unequal distribution, less about poverty than instability. This is often quoted to prove that ancient Chinese thinkers advocated egalitarianism. But honesty requires acknowledging that commentators have long disputed what Confucius meant by jun, translated here as evenness. Some read equal wealth; others read everyone securely occupying their proper place and receiving their appropriate share. The latter means each person's due, not everyone receiving the same. Both answers inhabit one classical sentence. The dispute has never been settled; successive generations have continued it.

The third takes us to Mesopotamia. Around thirty-seven hundred years ago, Babylon's Code of Hammurabi was carved on a stone pillar. A famous provision states, broadly, that if one free man blinds another, his own eye must be blinded in recompense. This is an eye for an eye, one of corrective justice's oldest forms. Its simple, powerful principle matches correction to harm, allowing neither leniency nor doubled revenge. But examine the qualifier free man. The same code prescribes completely different penalties for harming people of different ranks. Everyone being alike does not hold on that pillar: it is explicitly rejected. The stone mirrors how long humanity took to reach the apparently commonplace idea of equality before the law.

These sources are here not because the ancients were right but to establish two points. First, Friday afternoon's classroom argument concerns one of humanity's oldest documented public problems. Second, even the oldest traditions contain multiple voices. The Book of Rites favors need; this Analects sentence allows two readings; Hammurabi openly differentiates by status. Portraying any tradition as having always supported one exclusive form of justice is dishonest.

\section{Why One Rule Keeps Veering Off Course}
The class meeting produced no immediate decision. Ms. Li thought for two days, then announced a plan on Wednesday: the ten tablets would go to the fastest applicants. At eight that evening, students would add their names to a sign-up list in the class group chat, first come, first served. Her reasoning sounded unassailable: the same standard for all, no assessment of families, no marks, no special treatment. The procedure was as clean as distilled water.

The result was clear the next day. More than half the successful applicants already had computers, fast internet, and parents waiting by their phones at eight sharp. Lin Wan received nothing. Her mother's phone was charging that night, and her mother was on the night shift. Dachuan got one, though his family did not need another screen.

Pause at this outcome. It demonstrates the chapter's hardest phenomenon: \textbf{equal treatment can sometimes perpetuate actual inequality unchanged}. The rule addresses everyone in the same words, but their ability to respond differs completely. More than a century ago, French writer Anatole France mocked this majestic equality in The Red Lily: broadly, the law's sublime equality forbids rich and poor alike to sleep under bridges, beg in the streets, and steal bread. Literally, everyone is prohibited. In practice, the prohibition matters to only one kind of person.

Now hear four explanations of the same fastest-registration rule. These are not four emotions but four named traditions of justice, each with complete arguments and limitations.

\textbf{Explanation one: a clean procedure makes the outcome just, or libertarianism.}

A contemporary representative is American philosopher Robert Nozick. Its central claim, broadly, is that justice examines not the shape of wealth distribution but how each step occurred. If initial acquisition was legitimate and every transfer voluntary, without deception or coercion, an outcome where some have more and others less is just, however enormous the difference. Any distribution aiming to keep everyone approximately equal requires a continually intervening hand. Each intervention violates people's right to dispose freely of their belongings. From this position, fastest registration is unproblematic: the rule was public, nobody cheated, and justice follows automatically. Its strength is the forceful question of what entitles the state to touch my things. Its weakness is treating starting points as natural. Why was Lin Wan's mother's phone charging, and why was she working nights? The declaration that the procedure was lawful leaves these questions outside.

\textbf{Explanation two: ignore procedure and inspect the total, or utilitarianism.}

Founded by British thinkers including Bentham and Mill in the eighteenth and nineteenth centuries, utilitarianism broadly judges actions by the difference between the total happiness and total suffering they produce: create the greatest happiness for the greatest number. Applied to tablets, it changes the question. Ask not who deserves one but whose use produces the greatest educational benefit. Giving Dachuan a fourth screen adds almost nothing; giving Lin Wan one moves her from unclear lessons to clear ones. Utilitarianism would probably favor those in need, but for a reason unlike Lin Wan's own: not because she deserves it, but because this produces the best total. The position is practical, attentive to real consequences, and insists that every rule answer whether it works. It has two weaknesses. Happiness is difficult to aggregate and compare. More dangerously, maximizing the total can sacrifice minorities. Confiscating an innocent person's property might please a hundred people and yield a positive total, yet almost everyone would sense something wrong.

\textbf{Explanation three: the gap itself is the problem, or egalitarianism.}

Egalitarians answer that when a child's resources depend on whether her parents work nights, the gap itself is unjust; nothing further need be proved. The causes of the gap, such as family circumstances, talents, and birthplace, were not chosen by the child, who should not bear their moral cost. This position's strength is firmly removing luck from the moral ledger: you may answer for choices, not for the circumstances of your birth. Its weakness is practical too. If every difference disappears, is effort still rewarded? Dachuan's question about who would then stay up late is not purely selfish grumbling. Egalitarianism must answer rather than pretend it does not exist.

\textbf{Explanation four: ask not what people receive but what they can do with it, or the capability approach.}

Economist and philosopher Amartya Sen proposed a change of perspective. Distributive justice keeps asking who received how many things, yet identical things produce completely different lives in different hands. A tablet is a learning tool in a home with Wi-Fi and half a brick in one without it. Ask instead what each person is actually \textbf{able} to do and become. Can they see online lessons clearly, eat enough, and appear before classmates without shame? Sen studied famines in Bengal and elsewhere. His conclusion, broadly, was that many great famines occurred without inadequate food supplies. Rice filled shops, but poor people lost the entitlement and ability to obtain it and starved outside. This changed the world's understanding of poverty: it means not only lacking but being unable. The approach's strength is its unwavering attention to real life rather than ledger figures. Its difficulty is who determines which capabilities people should possess, returning us to competing rulers.

The four explanations are now before us. None can encompass the entire problem alone, and none is nonsense. This is not compromise for its own sake but a methodological warning: when a theory claims to explain everything, it has usually excluded some facts beforehand.

\section[Behind the Veil of Ignorance]{Further Challenge: Stand Behind the Veil of Ignorance}
Now introduce the chapter's weightiest intellectual tool. First, we must identify a shared ancestor of the preceding positions. From seventeenth-century Hobbes and Locke to eighteenth-century Rousseau, a European tradition called \textbf{social contract theory} broadly argued that legitimate rules are those people hypothetically agree upon together as equals. Governmental power does not fall from Heaven; it comes from a contract people make for security and order. This tradition influenced nearly all subsequent modern political debate. In A Theory of Justice (1971), American philosopher John Rawls brought it to a particularly sharp point through a thought experiment called the \textbf{veil of ignorance}.

Imagine people gathering to establish rules of distribution for the society they will inhabit: education, healthcare, opportunities. One condition applies. Until discussion ends, everyone stands behind a curtain concealing all information about themselves. You do not know your sex, family background, intelligence, health, appearance, or religion. You know only that when the curtain falls you will randomly become one member of this society and spend your life under the rules you have chosen.

Rawls argues that behind this curtain nobody would sign a contract oppressing the weak, because the gamble is unaffordable. You do not know whether you will be Lin Wan, the child whose mother works nights and whose phone is charging. Since you might draw the worst lot, rationality says to secure a safety net for those worst placed before discussing anything else. He derives two famous principles, broadly stated. First, everyone equally enjoys the most basic liberties of speech, belief, voting, and personal freedom. These cannot be discounted or traded: freedom cannot be exchanged for bread. Second, economic and status inequalities may be allowed only under two conditions: all positions are fairly open to everyone, and inequality leaves the worst off better placed than complete equality would.

Pause over the second principle. It means neither everyone gets the same nor everyone relies on their own abilities. It says, more intricately, that inequality is legitimate only when it serves the least advantaged. Consider high doctors' salaries. If they attract enough able people to medicine that even the poorest patients obtain care, the inequality may be just. If high salaries only speed wealthy patients' treatment while lengthening poor patients' queues, they fail the test. Rawls does not demand the elimination of differences. He demands that differences produce proof that they benefit those in the worst position.

Now bring the curtain into the classroom. Set tablet rules without knowing whether you are Lin Wan or Dachuan. Would you sign up to fastest registration? Probably not. You would more likely agree that those most in need receive protection, everyone has opportunities, and assessment does not publicly humiliate anyone. This is how the Rawlsian device works: it cannot decide for you, but it removes your own seat beneath you and makes you argue from reasons alone.

This tool also deserves substantial criticism. Libertarians ask why a curtain assuming universal risk aversion should decide for real people willing to gamble. Sen asks why, instead of imagining behind a curtain, we do not directly inspect what actual people can do. Their needed capabilities are precisely what the curtain conceals. Communitarians ask an even harder question: a person who knows neither themselves nor family, homeland, or faith does not exist. People always live through particular relationships and attachments. Can justice derived from a nonexistent person still be human justice?

Finally, an easily misread counterexample guards against a shallow reading. Many first conclude that Rawls therefore wants everyone to divide equally. Wrong: the veil permits inequality, even providing an official route for it, subject to that certificate of benefit. Mistaking Rawls for an advocate of equal shares resembles equating distribution according to work with letting the most capable take everything. These are two directions of the same carelessness. An intellectual tool's greatest danger is not opposition but reduction to a slogan used to batter others.

\section{Your School Distributes Things Too}
This ancient question now goes on trial daily at school, on phones, and at city street corners. Here are three recognizable examples.

\textbf{First: should financial-aid recipients be publicly listed?} Many schools identify students in poverty and post their names to meet procedural-justice requirements, citing transparency, supervision, and fraud prevention. Yet every year students would rather not apply than have their names appear. Having the whole class know your family is poor feels like public exposure at fourteen. Some schools therefore change their practice: no lists, speeches, or hardship contests. After assessment, funds quietly enter meal cards. Students with unusually low spending may even receive a subsidy notification without being asked to explain. Which is more just? Advocates of disclosure hold procedural justice: secrecy creates loopholes and kills fairness through hidden dealings. Advocates of discretion hold distributive justice and dignity: assistance should feed people, not require surrendering self-respect first. This is no drama of good against evil. Two forms of justice clash under the same banner of fairness.

\textbf{Second: neighborhood admissions and advanced classes.} Attending the nearest school sounds neutral: everyone follows distance, without favorites. But house prices arrive first. Those who can buy beside a prestigious school purchase proximity as privilege. A rule identical in wording carries starting-point differences intact into outcomes, the school-district housing version of France's bridge satire. Inside schools, grouping by marks likewise places children with after-school study support and children helping at family market stalls on different tracks. Acknowledging this need not abolish examinations. Dachuan's argument remains: effort needs a response. But every rule deserves periodic questioning about whose experience transforms its apparent sameness into something else.

\textbf{Third: algorithms assigning everyone the same rule.} Food-delivery platforms apply identical late penalties and route countdowns to hundreds of thousands of riders. The rule speaks the same words to everyone, but the platform unilaterally determines time estimates, allowances for traffic lights, and elevator waits. Riders are absent, certainly not behind any veil. Our checklist identifies the issue: not a dispute over distributive standards, but problems with question two, who distributes, and question four, whether procedure is just. Rule-makers are not governed by the rule; those governed cannot make it. Emergency rooms offer a contrasting case. Triage nurses prioritize severity over arrival order, and nobody protests unfairness because everyone knows they might become the heart-attack patient when the curtain falls. Daily life already mixes all four rulers. We simply rarely stop to identify which we hold.

\section{Your Question: If You Set the Rules}
Return to the classroom.

The Friday meeting starts again. This time Ms. Li adds a condition: you set the rules, but must do so behind the veil of ignorance. You do not know whether you are Lin Wan or Dachuan, whether your home has Wi-Fi, whether your mother works that day, or even whether your marks, running speed, or popularity are strong. You know only that when the curtain falls you will be one of the forty-eight.

What rules will you set?

Before writing, weigh every side again. Lin Wan is right that the same tablet creates different lives in different hands, and prioritizing urgent need satisfies both the total account and conscience. Dachuan is right that groups offering no response to effort gradually stop receiving it; misused need-based allocation rewards pleading poverty rather than contribution. Xiaoman is right that assessment authority is power and public declarations of family poverty cause harm. Chen Guo is right that procedure distributes opportunities: quick speakers, well-connected students, and confident volunteers always get another chance beyond the silent. Rawls reminds you to protect the worst position before discussing rewards. Nozick reminds you that each intervention for protection must justify its authority. Sen reminds you to count not only tablets but what people can actually do with them. The two-thousand-year-old voice in the Book of Rites suggests that the best class may be one where nobody must describe their misfortune to obtain something.

These reasons cannot all receive full marks simultaneously. Which you select, which you sacrifice, and why: that is your answer, not mine.

I want to leave one final question beyond these ten tablets. As an adult you will encounter larger distributions: wages, housing, hospital beds, school places, opportunities. Would you want to be treated as a number among identical people, or as a particular person receiving their due? If the latter, whom would you comfortably trust with the ruler that decides what you receive?

Please give your reasons. The silent people at the back of the classroom are waiting too.
